%%%%%%%%%%%%%%%%%%%%%%%%%%%%%%%%%%%%%%%%%%%%%%%%%%%%%%%%%%%%
%% CHAPTER 8 -- Adapting a Giant Model
%%%%%%%%%%%%%%%%%%%%%%%%%%%%%%%%%%%%%%%%%%%%%%%%%%%%%%%%%%%%

\chapter{Adapting a Giant Model}
\label{ntg:ch8}

\scenelabel
\begin{scenestrip}
\sceneitem{The problem}
\sceneitem{How we got here}
\sceneitem{Where we stand}
\sceneitem{What goes wrong}
\end{scenestrip}

\paragraph{The problem.} Picture a small product team that has just downloaded an
open-weight 70-billion-parameter language model from one of
the major laboratories. The model writes Python, summarises
French literature, drafts polite refusals to requests for
medical advice, and produces plausible recipes for cassoulet.
What it cannot do is reliably speak in the company's tone of
voice, follow the company's terminology conventions, or honour
the regulatory constraints of the jurisdiction the company will
be deploying it in.

The team has neither the budget nor the hardware to retrain the
model from scratch. They have a handful of accelerator cards
and a directory of labelled examples particular to their
domain. The chapter that follows is about the methods that turn
that handful of cards and that small directory into a deployable
adaptation.

The picture is not a special case. Almost no organisation
trains a base model from scratch; the cost is too high and the
expertise too specialised. Almost every organisation that
deploys a Large Language Model adapts one. The economic logic
of the field is shaped by this fact. The frontier laboratories
sell or release base models. Everyone else builds on top of
them. The chapter is about the techniques --- some technical,
some operational --- that make that building possible at the
budget of an ordinary product team.

\paragraph{How we got here.} The original idea, which the field tried first and largely
abandoned, was \emph{full fine-tuning}: take the pretrained
model, continue training it on your domain-specific data, and
ship the result. Full fine-tuning has the virtue of conceptual
simplicity. It also has several practical problems. The cost
is enormous, because the entire model has to be in memory and
every parameter has to be updated. The result is a full-sized
checkpoint per task, which is expensive to store and serve.
And the procedure has a tendency to overwrite useful general-
purpose knowledge with task-specific noise --- a phenomenon
called \emph{catastrophic forgetting}, after the visual image
of carefully accumulated capability being wiped out by a brief
adaptation pass.

The first response was \emph{adapters}. Neil Houlsby and his
collaborators at Google, in 2019, proposed inserting small
trainable layers --- so-called bottleneck layers --- into a
pretrained Transformer and freezing everything else. The
adapter parameters were a small fraction of the total but were
enough to recover most of what full fine-tuning would have
achieved. The adapter idea was good but architecturally
invasive: it required modifying the model's structure to make
room for the new layers.

The decisive idea came in 2021 from Edward Hu and his
collaborators at Microsoft. They called it \emph{Low-Rank
Adaptation}, or LoRA. The technical content rests on an
empirical observation that has held up well across many
settings. When a pretrained model is fine-tuned on a downstream
task, the changes to the model's weights are concentrated along
a small number of ``directions'' in weight space --- far fewer
than the model's parameter count would suggest. The technical
name for this is \emph{intrinsic dimensionality}, and the
empirical claim is that intrinsic dimensionality is small. LoRA
exploits this by parameterising the change to each weight
matrix as a product of two much smaller matrices, with the
total number of trainable parameters being a tiny fraction of
the original. The original weights are frozen; only the small
matrices are trained.

LoRA had two virtues that made it the field's default within
about eighteen months. The first was efficiency: a LoRA
adaptation of a 70-billion-parameter model uses on the order of
ten million trainable parameters --- about one ten-thousandth
of the full count. The second was modularity: a LoRA adapter
is a small file (tens of megabytes) that can be loaded onto
the base model at serve time, without touching the base
weights. An organisation that has trained a hundred different
adapters for a hundred different customers can serve all of
them from a single base model in memory, switching between
adapters per request. This is a clean operational picture, and
it is the production default for most organisations adapting
open models today.

The other thread came from below. As models grew, simply
\emph{loading} them onto an accelerator card became a problem.
A 70-billion-parameter model in standard 16-bit precision is
about 140 gigabytes; a single accelerator card typically has
40 to 80 gigabytes of memory. The remedy was \emph{quantisation}:
store the parameters in fewer bits, with some loss of precision.
Tim Dettmers and his collaborators at the University of
Washington showed in 2022, in a paper they called LLM.int8(),
that with careful per-channel scaling and outlier handling an
LLM could be served in 8-bit integers with negligible quality
loss. Subsequent methods --- GPTQ, AWQ, SmoothQuant --- pushed
the precision lower while preserving quality. By 2023 a
70-billion-parameter model could be served in 4 bits at
roughly 35 gigabytes, which fits on the cards a small team can
plausibly afford.

In 2023 Dettmers's team unified the two threads in a method
called \emph{QLoRA}. The recipe is to quantise the base model
to 4 bits, freeze it, and train a LoRA adapter on top of it in
higher precision. The combination made the adaptation of
70-billion-parameter models tractable on a single workstation-
class graphics card. It is the production default for most
adaptation work today.

\paragraph{Where we stand.} The modern view treats a deployed model not as a single
monolithic artefact but as a pair: a large, rarely-changing
base model stored in compressed form, and a small, frequently-
changing adaptation stored separately. The base is shared
across tasks and customers. The adaptation is what makes any
particular checkpoint specific. The adaptation can be a LoRA
adapter, a fine-tuned subset of the model's parameters, or
sometimes just a carefully designed prompt template.

Engineers at organisations that deploy adapted models manage
their adapters the way they once managed model checkpoints --- a
registry of versioned artefacts, with metadata about which base
model they were trained against, which dataset, and which
evaluation suite they passed. A serving system might host a
dozen LoRA adapters in front of a single base, with a router
selecting among them per request based on which customer is
calling.

Quantisation formats now include 16-bit, 8-bit, and 4-bit
variants, each with several flavours optimised for different
properties of the model. Newer formats explicitly designed for
LLM serving --- the so-called NF4, for instance --- are
information-theoretically motivated to preserve as much of the
model's behaviour as possible at low bit depth. The choices in
this area have stabilised enough that picking the right
quantisation for a given deployment is more an engineering
question than a research one.

\paragraph{What goes wrong.} The mathematical assumption behind LoRA --- that the directions
of useful adaptation are few --- is empirical. For some tasks
the assumption holds well and a small adapter recovers nearly
all of full fine-tuning. For other tasks, especially those that
require the model to learn substantially new behaviour rather
than to refine existing behaviour, the small adapter
underperforms in ways that may not show up on the benchmarks
the team chose. Diagnosing this requires evaluation that goes
beyond the headline numbers, which not every team does.

Quantisation introduces rounding errors that are, in
expectation, small but that can be decisive on rare or
unusual inputs. A model that scores within a small margin of
its full-precision counterpart on a standard benchmark may
give noticeably worse answers on the inputs that lie outside
the benchmark's distribution. The qualitative pattern is that
quantisation errors disproportionately affect the model's
behaviour on edge cases --- which is to say, on exactly the
inputs where careful behaviour matters most.

There is an operational class of failure that has nothing to
do with the math. An adapter trained on out-of-date documents
will produce out-of-date answers, gracefully and confidently.
An adapter trained without a regression test suite will
silently regress capabilities the base model had. An adapter
trained on a dataset that was, unbeknownst to the team,
contaminated with the test set will appear to perform much
better than it does. None of these failures are visible from
inspecting the adapter file. All of them are visible from
careful evaluation, which is why the engineering practice
around adapter management is at least as important as the
mathematical apparatus that makes the adapters cheap.

\section{The idea, in plain language}
\addcontentsline{toc}{section}{The idea, in plain language}

\subsection*{LoRA as the editor's correction}

The most useful image for understanding LoRA is the editor's
correction to a manuscript. The base model is the manuscript:
hundreds of pages of carefully written prose, the work of
months. Full fine-tuning would be rewriting the manuscript from
scratch. LoRA is what an editor does: a few hundred small
margin notes that, applied at the right places, change the
manuscript's tone, its terminology, its emphasis, without
touching most of the text.

The image is faithful in its key features. The base model is
unchanged after LoRA training; the adapter is a small set of
corrections that overlay it. The corrections can be applied or
removed at will, just as an editor's notes can be accepted or
rejected. Multiple sets of corrections can be maintained
separately --- one editor's notes for the legal version, another
editor's notes for the marketing version, a third editor's
notes for the version going to the regulator --- and the same
underlying manuscript can be served with any of them, depending
on who is asking.

The image is also misleading in two ways worth noting. First,
the corrections in LoRA are not localised to particular
passages of the model the way an editor's notes are localised
to particular passages of a manuscript. The corrections are
distributed across all the model's weights, in a structured
mathematical form that makes them small in total size but
diffuse in influence. Second, the model is not a manuscript
that humans read; it is a procedure that produces output, and
the corrections affect the procedure. ``Corrections at the
right places'' is a metaphor for ``corrections in the right
mathematical directions'', which is not a thing the human
intuition for editing has a counterpart for.

\subsection*{Why the low-rank bet is usually a good one}

The editor's correction analogy tells you what LoRA does. It leaves
open the harder question: why should restricting those corrections
to a compact form work at all? Why would the thing you want to
change fit into a small corner of the available change-space?

The answer begins with a close look at what typical adaptation
tasks actually require. A medical institution wants the model to
use clinical vocabulary, hedge claims about patient outcomes, and
decline to offer unauthorised treatment recommendations. A law firm
wants it to follow a particular citation style and flag disputed
facts for counsel. A customer-service department wants a particular
tone and a particular escalation threshold.

These requirements have something in common: they are changes to
the model's behavioural \emph{defaults}, not demands that it learn
new knowledge it has never encountered. The model already knows
what clinical vocabulary is. It already knows what a hedged
statement looks like. It already knows several styles of refusal.
What the adaptation is doing is tilting the model's choices
among things it already knows how to produce --- shifting its
defaults along a small number of behavioural axes.

Think of it this way. A 70-billion-parameter model has, in
principle, a vast space of possible behaviours --- an almost
unimaginably large number of ways its outputs could be different
from what they currently are. Almost none of those possible
differences correspond to what a real adaptation task is asking
for. Shifting from ``general assistant'' to ``assistant that
prefers clinical vocabulary'' requires changing something
relatively narrow and specific. The rest of the enormous
behaviour-space --- the model's grasp of chemistry, its knowledge
of 19th-century history, its ability to write a sonnet --- can
stay exactly as it is.

LoRA exploits this by restricting the adaptation to exactly the
kind of change that adaptation tasks actually require: a small
number of adjustments to the model's default tendencies, encoded
compactly as a product of two small matrices. Because the useful
change is narrow, the compact encoding is adequate. The redundant
99.9\% of the parameter space is left untouched, which is also
why catastrophic forgetting is reduced: if you do not touch the
parameters that encode chemistry and history, you do not degrade
the model's knowledge of chemistry and history.

The bet fails when the task requires the model to acquire
genuinely new capabilities that were absent from pretraining ---
not to shift its defaults among things it already knows, but to
learn something it has never encountered. A LoRA adapter trained
on a new symbolic reasoning framework the model has never seen
will underperform a model that was pretrained on that framework
from the beginning. The compact encoding cannot invent capability
from nothing; it can only redirect capability that already exists.
Knowing this is the engineering judgment that decides whether LoRA
is the right tool or whether a different approach is needed.

\subsection*{Quantisation as the photograph compressed to a JPEG}

The most useful image for understanding quantisation is the
photograph compressed to a JPEG file. A high-resolution
photograph stored in raw format is large but precise: every
pixel's colour is recorded with full precision. The same
photograph saved as a JPEG is much smaller, because the JPEG
algorithm replaces the precise values with approximations
chosen to be visually close. The JPEG is, on the whole,
indistinguishable from the original to a casual viewer. On
careful inspection it has visible artefacts: jagged edges
around sharp transitions, slight banding in smooth gradients,
loss of fine texture in high-frequency areas.

Quantising a Large Language Model is the same trick. The
parameters of the model are stored, after quantisation, in
fewer bits per parameter, with some loss of precision. The
loss of precision is, on the whole, indistinguishable in the
model's typical behaviour. On careful inspection it has visible
artefacts: degraded performance on rare inputs, reduced
calibration, occasional surprising answers in places where the
unquantised model would have done better. The JPEG analogy
even extends to the trade-off curve: more aggressive
compression saves more space but produces more visible
artefacts, and there is a sweet spot for any given application.

The cleverness of the recent quantisation work is in the
choice of which precision to throw away where. A naive
implementation that simply rounded all parameters to four bits
would produce a model that did not work. The methods that work
well --- LLM.int8(), GPTQ, AWQ, NF4, and others --- preserve
precision in the parameters that matter most, accept loss in
the parameters that matter less, and handle the few outlier
parameters that have to be kept in higher precision separately.
The result is a 4-bit model whose behaviour is much closer to
the unquantised baseline than naive 4-bit rounding would
suggest.

\subsection*{When adaptation is appropriate, and when it is not}

Several situations call for adaptation. A company that wants
the model to use its terminology consistently --- replacing
``client'' with ``patient'' in a medical setting, for instance,
or honouring a particular product naming convention --- can
adapt a model on a corpus that uses the desired terminology. A
company that wants the model to refuse requests outside its
domain, with a particular wording, can adapt a model on
examples of the desired refusal. A company that wants the model
to follow a particular response format --- always ending with a
disclaimer, always citing sources, always producing structured
JSON --- can adapt a model on examples of the desired format.

Several situations do not call for adaptation. A company that
wants the model to know its product manual is usually better
served by retrieval (Chapter~10) than by adaptation: retrieval
gives the model fresh access to the document, where adaptation
would freeze the document into the weights and require
retraining whenever the document changes. A company that wants
the model to know who the current CEO is should not put that
into the weights at all; it should put it into the prompt or
into a tool the model can call. The general rule is that
adaptation is appropriate for stable patterns of behaviour and
inappropriate for rapidly-changing facts.

\section{The engineering consequence}
\addcontentsline{toc}{section}{The engineering consequence}

\paragraph{The base model and the adapter are separable
artefacts.} A team that has trained an adapter against one
version of a base model usually has to retrain it against the
next version. Vendors sometimes change the base model they
serve under the same API name --- ``GPT-4'' has, over the past
two years, referred to several different underlying models with
different post-training. An adaptation built against one
version of a vendor-served model can drift in behaviour when
the vendor updates the underlying model. This is a real
operational concern for any team building a long-lived product
on top of a vendor-served model.

\paragraph{Adapter composition is not what it appears.} An
adapter for tone of voice and an adapter for terminology
cannot, in general, be combined by addition: the resulting
weights do not faithfully represent the union of behaviours.
There is research on adapter composition; there is no clean
recipe that works in all cases. A team that needs multiple
behaviours from one model is usually better served by training
a single adapter on a combined dataset than by training
specialised adapters separately and combining them at serve
time.

\paragraph{Quantisation degrades behaviour at the tails, not
in the middle.} Standard benchmarks usually evaluate models on
distributions close to the training distribution. Quantisation
errors are, in expectation, small on such distributions. They
are, in expectation, larger on rare or out-of-distribution
inputs. A team evaluating a quantised model should test it on
the distribution it will actually face in deployment, not on
the standard benchmarks alone. The quantisation budget is one
of the more under-evaluated decisions in deployment.

\paragraph{The cost story is now in serving, not in training.}
For a team adapting an open base model, the upfront cost of
adaptation is small (a single workstation-class GPU for a few
hours to a few days). The ongoing cost is in serving: every
query consumes some amount of GPU time. The unit economics of
the resulting product depend on how aggressively the model can
be quantised, how efficiently it can be served, and how many
queries the model can process per second per GPU. The
infrastructure work of Chapter~7 is what determines those
numbers.

\section{What goes wrong, and why}
\addcontentsline{toc}{section}{What goes wrong, and why}

\paragraph{Catastrophic forgetting.} A model that is fine-tuned
on a narrow corpus can lose capability that was present in the
base model but not represented in the adaptation corpus. A
model that is adapted on legal documents can become measurably
worse at conversational helpfulness; a model adapted on code
can become worse at prose. LoRA is partly designed to mitigate
this --- by freezing the base weights, it limits the damage --- but
it does not eliminate the problem. The fix is regression
testing: maintaining a suite of evaluations that exercise
capabilities that are not the target of the adaptation, and
running them every time an adapter is updated.

\paragraph{Quiet quality regressions.} A team that ships a
quantised version of a model in place of an unquantised one,
without testing on the deployment distribution, can introduce
quality regressions that are invisible in the benchmark numbers
but visible to users. The regressions tend to be at the
edges --- on rare inputs, on long sequences, on tasks that
require precise reasoning. The fix is, again, evaluation on
the deployment distribution.

\paragraph{Adapter drift over base updates.} An adapter trained
against one version of a base model may behave subtly
differently when the base is updated. Vendors that update their
base models without notification put their customers in the
position of debugging changes they did not make and cannot
investigate. The fix, where possible, is to pin to a specific
base model version and to test before any update. The fix is
not always possible, because not all vendors expose version
pinning.

\paragraph{Confidentiality concerns in adaptation data.} The
data used to adapt a model becomes, in a precise sense,
encoded in the adapter. A model adapted on internal documents
can, under the right prompts, reveal the contents of those
documents. A team adapting a model on confidential data should
treat the resulting adapter with the same security posture as
the data itself --- not as a derivative work. Several of the
high-profile data-leak incidents in LLM products have been
adaptation-data leaks dressed up as model leaks.

\section*{Bridges}
\addcontentsline{toc}{section}{Bridges}

This chapter built on Chapter~3 (the deep neural network whose
parameters we are now adapting), Chapter~6 (the pretraining
procedure that produced the base model), and Chapter~7 (the
serving infrastructure within which adapters are deployed). It
described how organisations other than the frontier laboratories
turn open base models into products. Chapter~9 describes a
particular kind of adaptation --- the alignment work --- that the
frontier laboratories themselves do to turn a base model into
an assistant. Chapter~10 describes an alternative to adaptation
for many use cases, retrieval, which avoids freezing rapidly-
changing knowledge into the model. Chapter~12 returns to the
deployment of adapted models from the perspective of the
developer who is calling them through an API. The next chapter
--- \emph{Teaching Manners and the Cost of Doing So} --- explains
how the alignment that turned the GPT-3 base model into ChatGPT
actually works, and what it costs.
